\section{Formal Model: Task Lifecycle}\label{sec:formal-model}

The P-A-R runtime is designed around a task lifecycle that is dynamic in nature: agents may spawn tasks, suspend them awaiting external input, resume them, or cancel them. Without a formal specification, such a system is prone to illegal state transitions that manifest only at runtime, leading to crashes in production environments. To address this, we model the task lifecycle as a Labeled Transition System (LTS) that captures precisely which transitions are permitted from each state. This model serves as the authoritative reference for the type-safe state machine implemented in the compiler, and provides the basis for the formal properties we prove in \S\ref{sec:formal-properties}.

\subsection{State Space and Transitions}\label{sec:lts-definition}

\begin{definition}[State Space]\label{def:states}
The set of task states $\mathcal{S} = \{\,\identifiable{Pending}, \identifiable{Scheduled}, \identifiable{Running}, \identifiable{Waiting\_input}, \identifiable{Suspended}, \identifiable{Completed}, \identifiable{Failed}, \identifiable{Cancelled}\,\}$.
\end{definition}

\begin{definition}[Terminal States]\label{def:terminal-states}
The subset $\mathcal{T} = \{\,\identifiable{Completed}, \identifiable{Failed}, \identifiable{Cancelled}\,\}$ is the set of terminal states. No outgoing transitions exist from any $t \in \mathcal{T}$.
\end{definition}

\begin{definition}[Transition Labels]\label{def:labels}
The set of transition labels $\mathcal{L} = \{\,\tau,\, \identifiable{schedule},\, \identifiable{cancel},\, \identifiable{suspend},\, \identifiable{resume},\, \identifiable{wait},\, \identifiable{provide},\, \identifiable{fail}\,\}$.
\end{definition}

\begin{figure}[t]
  \centering
  \begin{tikzpicture}[
    node distance=2.2cm and 3cm,
    every edge/.style={draw, ->, >=stealth},
    state/.style={ellipse, draw, minimum width=2.2cm, minimum height=0.9cm},
    terminal/.style={ellipse, draw, double, double distance=2pt, minimum width=2.2cm, minimum height=0.9cm},
    pending/.style={state, fill=green!30},
    scheduled/.style={state, fill=blue!15},
    running/.style={state, fill=blue!15},
    waiting/.style={state, fill=blue!15},
    suspended/.style={state, fill=blue!15},
    completed/.style={terminal, fill=red!20},
    failed/.style={terminal, fill=red!20},
    cancelled/.style={terminal, fill=red!20},
    label/.style={midway, fill=white, font=\small}
  ]

  \node[pending] (pending) {$Pending$};
  \node[scheduled, right=of pending] (scheduled) {$Scheduled$};
  \node[running, right=of scheduled] (running) {$Running$};
  \node[waiting, below=of running] (waiting) {$Waiting\_input$};
  \node[suspended, below=of scheduled] (suspended) {$Suspended$};

  \node[completed, below right=2.5cm and 1cm of running] (completed) {$Completed$};
  \node[failed, below=of completed] (failed) {$Failed$};
  \node[cancelled, right=of completed] (cancelled) {$Cancelled$};

  \path
    (pending) edge node[label] {$schedule$} (scheduled)
    (pending) edge[bend left=30] node[label] {$cancel$} (cancelled)

    (scheduled) edge node[label] {$\tau$} (running)
    (scheduled) edge[bend right=20] node[label] {$cancel$} (cancelled)

    (running) edge node[label] {$wait$} (waiting)
    (running) edge node[label] {$suspend$} (suspended)
    (running) edge node[label] {$cancel$} (cancelled)
    (running) edge node[label] {$fail$} (failed)
    (running) edge node[label] {$\tau$} (completed)

    (waiting) edge node[label] {$provide$} (running)
    (waiting) edge node[label] {$cancel$} (cancelled)
    (waiting) edge node[label] {$fail$} (failed)
    (waiting) edge node[label] {$\tau$} (completed)

    (suspended) edge node[label] {$resume$} (scheduled)
    (suspended) edge[bend right=15] node[label] {$resume$} (running)
    (suspended) edge node[label] {$cancel$} (cancelled)
    (suspended) edge node[label] {$fail$} (failed)
    (suspended) edge node[label] {$\tau$} (completed)
  ;
  \end{tikzpicture}
  \caption{State transition diagram for the P-A-R task lifecycle. Terminal states are double-circled. Intermediate states are shown in blue, and the initial \identifiable{Pending} state in green.}
  \label{fig:state-machine}
\end{figure}

\subsection{Inference Rules}\label{sec:inference-rules}

We present the operational semantics of the task lifecycle as a set of inference rules. Each rule has the form $\inferrule[Label]{premise_1 \and ... \and premise_n}{conclusion}$, encoding the condition under which a transition may fire.

\begin{mathpar}
\inferrule[Pending-Schedule]
  {\identifiable{validate}(s, \identifiable{schedule}) = \identifiable{ok}}
  {s \xrightarrow{\identifiable{schedule}} s'}

\inferrule[Pending-Cancel]
  {\identifiable{validate}(s, \identifiable{cancel}) = \identifiable{ok}}
  {s \xrightarrow{\identifiable{cancel}} \identifiable{Cancelled}}

\inferrule[Scheduled-$\tau$]
  {\identifiable{validate}(s, \tau) = \identifiable{ok}}
  {s \xrightarrow{\tau} \identifiable{Running}}

\inferrule[Scheduled-Cancel]
  {\identifiable{validate}(s, \identifiable{cancel}) = \identifiable{ok}}
  {s \xrightarrow{\identifiable{cancel}} \identifiable{Cancelled}}

\inferrule[Running-Wait]
  {\identifiable{validate}(s, \identifiable{wait}) = \identifiable{ok}}
  {s \xrightarrow{\identifiable{wait}} \identifiable{Waiting\_input}}

\inferrule[Running-Suspend]
  {\identifiable{validate}(s, \identifiable{suspend}) = \identifiable{ok}}
  {s \xrightarrow{\identifiable{suspend}} \identifiable{Suspended}}

\inferrule[Running-Cancel]
  {\identifiable{validate}(s, \identifiable{cancel}) = \identifiable{ok}}
  {s \xrightarrow{\identifiable{cancel}} \identifiable{Cancelled}}

\inferrule[Running-Fail]
  {\identifiable{validate}(s, \identifiable{fail}) = \identifiable{ok}}
  {s \xrightarrow{\identifiable{fail}} \identifiable{Failed}}

\inferrule[Running-Complete]
  {\identifiable{validate}(s, \tau) = \identifiable{ok}}
  {s \xrightarrow{\tau} \identifiable{Completed}}

\inferrule[Waiting-Provide]
  {\identifiable{validate}(s, \identifiable{provide}) = \identifiable{ok}}
  {s \xrightarrow{\identifiable{provide}} \identifiable{Running}}

\inferrule[Waiting-Cancel]
  {\identifiable{validate}(s, \identifiable{cancel}) = \identifiable{ok}}
  {s \xrightarrow{\identifiable{cancel}} \identifiable{Cancelled}}

\inferrule[Waiting-Fail]
  {\identifiable{validate}(s, \identifiable{fail}) = \identifiable{ok}}
  {s \xrightarrow{\identifiable{fail}} \identifiable{Failed}}

\inferrule[Waiting-Complete]
  {\identifiable{validate}(s, \tau) = \identifiable{ok}}
  {s \xrightarrow{\tau} \identifiable{Completed}}

\inferrule[Suspended-Resume-Scheduled]
  {\identifiable{validate}(s, \identifiable{resume}) = \identifiable{ok}}
  {s \xrightarrow{\identifiable{resume}} \identifiable{Scheduled}}

\inferrule[Suspended-Resume-Running]
  {\identifiable{validate}(s, \identifiable{resume}) = \identifiable{ok}}
  {s \xrightarrow{\identifiable{resume}} \identifiable{Running}}

\inferrule[Suspended-Cancel]
  {\identifiable{validate}(s, \identifiable{cancel}) = \identifiable{ok}}
  {s \xrightarrow{\identifiable{cancel}} \identifiable{Cancelled}}

\inferrule[Suspended-Fail]
  {\identifiable{validate}(s, \identifiable{fail}) = \identifiable{ok}}
  {s \xrightarrow{\identifiable{fail}} \identifiable{Failed}}

\inferrule[Suspended-Complete]
  {\identifiable{validate}(s, \tau) = \identifiable{ok}}
  {s \xrightarrow{\tau} \identifiable{Completed}}
\end{mathpar}

Crucially, no inference rules are defined for transitions emanating from terminal states, enforcing the constraint that once a task reaches \identifiable{Completed}, \identifiable{Failed}, or \identifiable{Cancelled}, it cannot transition further. This property is stated formally below.

\subsection{Formal Properties}\label{sec:formal-properties}

We prove the following invariants of the LTS. All seven properties are validated empirically in \S\ref{sec:evaluation}.

\begin{prop}[Terminality]\label{prop:terminality}
For every terminal state $t \in \mathcal{T}$, there exists no label $l \in \mathcal{L}$ such that $t \xrightarrow{l} s$ for any state $s$ (validated empirically in \S\ref{sec:evaluation}).
\end{prop}

\begin{prop}[No Self-Transitions]\label{prop:no-self}
For all states $s \in \mathcal{S}$ and labels $l \in \mathcal{L}$, the transition $s \xrightarrow{l} s$ does not occur (validated empirically in \S\ref{sec:evaluation}).
\end{prop}

\begin{prop}[Deterministic Guards]\label{prop:deterministic-guards}
The guard function $\identifiable{validate} : \mathcal{S} \times \mathcal{L} \to \{\identifiable{ok}, \identifiable{error}\}$ is total and deterministic: for every $(s, l)$ pair, $\identifiable{validate}(s, l)$ returns exactly one outcome (validated empirically in \S\ref{sec:evaluation}).
\end{prop}

\begin{prop}[Maximal Execution Path]\label{prop:max-path}
From the initial state \identifiable{Pending}, any valid execution path to a terminal state contains at most 7 transitions (validated empirically in \S\ref{sec:evaluation}).
\end{prop}

\begin{prop}[Max Path Proof Sketch]\label{prop:max-path-proof}
Consider the longest feasible path. From \identifiable{Pending}, the maximal sequence is:

\[
\begin{aligned}
\identifiable{Pending}
& \xrightarrow{\identifiable{schedule}} \identifiable{Scheduled} \\
& \xrightarrow{\tau} \identifiable{Running} \\
& \xrightarrow{\identifiable{suspend}} \identifiable{Suspended} \\
& \xrightarrow{\identifiable{resume}} \identifiable{Scheduled} \\
& \xrightarrow{\tau} \identifiable{Running} \\
& \xrightarrow{\tau} \identifiable{Completed}
\end{aligned}
\]

This trace contains exactly 7 transitions, and no valid path can exceed this length because each state along the path is visited at most once.
\end{prop}

\begin{prop}[Type Safety]\label{prop:type-safety}
The OCaml variant type $\identifiable{handler\_result}$ enumerates all five non-terminal transitions (\identifiable{schedule}, \identifiable{cancel}, \identifiable{suspend}, \identifiable{resume}, \identifiable{wait}, \identifiable{provide}, \identifiable{fail}) plus the silent $\tau$ label. Pattern matching on this type is exhaustive by construction (validated empirically in \S\ref{sec:evaluation}).
\end{prop}

\begin{prop}[Exhaustiveness]\label{prop:exhaustiveness}
The OCaml compiler emits a warning when pattern matching on $\identifiable{handler\_result}$ does not cover all constructors. This warning is elevated to a compile-time error in P-A-R, forcing the programmer to handle every possible transition (validated empirically in \S\ref{sec:evaluation}).
\end{prop}

\begin{prop}[Semantic Validity]\label{prop:semantic-validity}
The guard function $\identifiable{validate}$ must be called on every transition attempt. An implementation that bypasses $\identifiable{validate}$ and directly applies a transition may violate any of the preceding properties (validated empirically in \S\ref{sec:evaluation}).
\end{prop}

\subsection{Motivating Example}\label{sec:motivating-example}

To illustrate why a type-safe state machine matters in practice, consider a handler that updates task status. In Python, one might write:

\begin{verbatim}
def complete_task(t):
    if t.status == "running":
        t.status = "completed"
    else:
        raise InvalidStateError(...)
\end{verbatim}

This code compiles without error, but at runtime it may encounter a task in state \identifiable{Waiting\_input} that is also considered valid for completion. The guard logic is implicit, scattered across conditionals, and easy to violate. Worse, adding a new state (e.g., \identifiable{Suspended}) requires hunting through all such conditionals to update the logic.

In OCaml, the same domain is captured by the $\identifiable{handler\_result}$ variant:

\begin{verbatim}
type handler_result =
  | Scheduled of task
  | Waiting_input of task
  | Completed of task
  | Failed of task
  | Cancelled of task
  | InvalidTransition of task * transition
\end{verbatim}

Pattern matching on a value of this type forces the programmer to handle all five non-terminal states explicitly. A missing case triggers a compiler warning, which P-A-R elevates to a hard error. The type system thus makes illegal state transitions a compile-time failure rather than a runtime crash. This contrast is representative: dynamically typed frameworks defer correctness checks to runtime, whereas P-A-R makes them inescapable at compile time.